% !TeX root = ../papers/nc-p3.tex
\section{Kummer曲面の極写像と十六の二重接触平面}
\label{sec:kummer-polar-completion}
前節の通常倍率を保ち、$A_0,A_1$を成分順$(00,01,10,11)$で列ベクトルと読む。
$u=z(\lambda^2+1)$, $v=\lambda(z+1)$, $w=\lambda^2z^2+1$として
\[
 T=\begin{pmatrix}
 u&v&zv&w\\v&\lambda^2+1&w&v\\zv&w&zu&zv\\w&v&zv&u
 \end{pmatrix}.
\]
\begin{theorem}\label{thm:kummer-polar-completion}
$F=0$の有理写像として$[A_1]=[T\nabla F(A_0)]$である。
この写像は対合で、窓の二位置移動を文字ラベルと対角倍率の変更まで実現する。
従って$T$は曲面の射影双対を再び同じ曲面に同定する。
十六の二重接触平面の各切断は滑らかな二次曲線を二重に数えたものである。
これらの上の許容窓は全て非標準ASである。
\end{theorem}
\begin{proof}
$pq\ne0$で$x=-m/(2p),y=-k/(2q)$を代入すると、完成行列の四成分は
\[
 \left(\frac{1+e_0(h_1-e_1)}{2h_1}, \frac{f_0g_1-g_0f_1-1}{2g_1}, \frac{g_0f_1-f_0g_1-1}{2f_1}, \frac{1+e_0(e_1-h_1)}{2e_1}\right)
\]
である。$\det T=(z^2-1)^2(\lambda^2z-1)^4$。
\path{computations/verify_kummer_polar.py}は各成分の分母を払い、
上の列と$T\nabla F$の六つの二次小行列が全て
$\mathbb Q[a,b,c,d,z,\lambda]$で$F$で割り切れることを検算する。
従って滑らかな点で両方の表示が定まる限り射影値は等しい。

$A_0,A_1$を交換すると$C,D$も交換され、両者のSchur補を対角と非対角に
分けた四つの条件がそのまま保たれる。テンソルについて具体的に、
元の成分$w_{k,l\mathbin\oplus1,i,j\mathbin\oplus1}$へ
$q^{I+J-K-L}$を掛ける。ここで$q^2=\lambda$、大文字は各添字が
$\{2,3\}$に属すとき1、他で0である。薄い色では指数が零で二つの薄い係数の積が
元の規格化に戻る。厚い色では$I=J,K=L$で、左上・右下は不変、
右上は$1/\lambda$、左下は$\lambda$倍となる。
従って得られるテンソルは同じ$(z,\lambda)$で$A_0,A_1$を交換したものに等しい。
この関係は二重根以外にも成立する。これは二位置移動を伴う関係であり、
移動を忘れた添字保持同型を自動的には与えない。

同コードで$F(T\nabla F)$も$F$で割り切れることを確認する。
どの許容倍率でもGauss像は二次元である。実際、通常二重点での接平面の極限は、
その非退化な二次接錐の各母線での接平面を全て含む。これは節点を一度吹き上げ、
例外二次曲線で勾配の先頭項を読むことで分かる。一節点につき双対平面内に
滑らかな二次曲線が一つあり、異なる節点ではその平面が異なる。
既約なGauss像の閉包はこれら二曲線を含むため次元は二である。
よって$T$を施した像は既約な$F=0$に稠密であり、両行列とも$pqD_0\ne0$の
一般点で交換できる。重根による完成が一意なので、再度の交換で元の行列に戻る。
この同定は射影双対性の双対定理と$T$の対称性からも従う。

十六平面を列挙する。四平面は
$b=c,\ c=z^2 b,\ d=za,\ a=zd$であり、
四平面は$f_1,g_1,e_1,h_1$の分子を零とするものである。残る八平面は
\[
 b=\lambda d,\quad d=\lambda z b,\quad c=\lambda z d,\quad d=\lambda c,
 \qquad
 a=\lambda z b,\quad b=\lambda a,\quad a=\lambda c,\quad c=\lambda z a .
\]
各法線を$L$と書く。同コードはチャートで単元の係数を一つ選んでその座標を消去し、
$F$の制限が単元倍の二次式の平方で、二次式のHesse行列式が単元であることを
因数分解で調べる。それぞれの二次曲線上の滑らかな点の接平面は$L$である。
$TL$は元の曲面の節点である。八つの交差成分平面では、$TL$を
$a=1$に規格化した値は定理\ref{thm:klein-all-pairings}の表の八つの$A_0$と一致する。
$d\in\{z,1/z\}$、$c/b\in\{1,z^2\}$、各組に二つの相異なる$b$があり、
$\lambda^2z\ne1$によりいずれも可逆である。他の八つの$TL$は四座標頂点と
\[
 (1,\lambda,1/\lambda,1),\quad(1,1/(\lambda z),\lambda z,1),\quad
 (1,\lambda,\lambda z,\lambda^2z),\quad
 (1,1/(\lambda z),1/\lambda,1/(\lambda^2z))
\]
で、全て$D_0=0$である。これらも通常倍率で相異なる。
これで前節の全十六節点と十六平面が対をなす。
ほかに平面切断が二重の滑らかな二次曲線となる平面はない。
実際、Gauss対合がその曲線を一点に縮めれば、逆写像はその点で
正則に定まらず、$T$が可逆なのでそこは節点である。

最初の四平面では$n=0$または$o=0$であり、真の二重根なら$x=0$または$y=0$に
なって可逆チャートを外れる。次の四平面では$pq=0$で、重判別式と有限な解を
仮定すると対応する二次式は恒等的に零になる。これらは前二節で処理した。
最後の八平面では$A_1$を新たな基底とすれば$P\equiv Q\equiv0$である。
実際、平面接触により一般点の$[A_1]$は上の特定の$[TL]$であり、
$pq\ne0$で有限な重根表示は正則なので節点への特殊化にも同じ値が保たれる。
従って交換した窓を定理\ref{thm:all-zero-pairing-defect}で除外できる。
二位置移動と基底変更は正則性を保つから元の窓も非正則である。
\end{proof}
